\documentclass[tikz, border=10pt]{standalone}
\usepackage{tikz}
\usepackage{amsmath}
\usepackage{amssymb}
\usepackage{xcolor}
\usetikzlibrary{arrows.meta, decorations.markings, angles, quotes, calc}

\definecolor{waveBlue}{HTML}{1a6eb5}
\definecolor{sensorBlue}{HTML}{2196f3}
\definecolor{pathRed}{HTML}{d32f2f}
\definecolor{annoOrange}{HTML}{e65100}

\begin{document}
\begin{tikzpicture}[
    font=\sffamily\small,
    >=Stealth
]

% ── Parameters (all in TikZ units = cm) ─────────────────────────────────
% 8 sensors, spacing d represented as 1.8 units
\def\d{1.8}       % element spacing in drawing units
\def\N{8}         % number of elements shown
\def\theta{35}    % bearing from broadside in degrees → wavefront angle from vertical

% ── Array baseline ───────────────────────────────────────────────────────
\foreach \i in {0,1,...,7}{
    \pgfmathsetmacro{\x}{\i * \d}
    \draw[sensorBlue, fill=sensorBlue!20, thick] (\x, 0) circle (0.2);
    \node[below=4pt, font=\sffamily\footnotesize] at (\x, 0)
        {$n{=}\i$};
}
% Array axis label
\node[below=18pt] at (3.5*\d, 0) {ULA axis};

% ── Element spacing annotation ───────────────────────────────────────────
\draw[<->, thin] (0, -0.7) -- (\d, -0.7)
    node[midway, below, font=\sffamily\footnotesize]
    {$d = 0.1875\,\text{m}$};

% ── Incoming wavefront (diagonal line) ──────────────────────────────────
% Wavefront arrives from upper-left; the wavefront is perpendicular to
% the direction of propagation. Direction of propagation makes angle
% (90° - theta) with the array axis.
% Wavefront line passes through element 0 at y=0.
% Wavefront direction vector: (sin(theta), -cos(theta)) normalised
% We draw the wavefront as a line with slope = -cot(theta) shifted right

\pgfmathsetmacro{\waveslope}{cos(\theta)/sin(\theta)}   % slope dy/dx of wavefront

% Pre-compute start/end y-coordinates (TikZ cannot evaluate \macro*expr inline)
\pgfmathsetmacro{\wfStartY}{\waveslope*(-1.0)}
\pgfmathsetmacro{\wfEndX}{7*\d + 0.5}
\pgfmathsetmacro{\wfEndY}{\waveslope*(7*\d + 0.5)}

% Wavefront line through element 0 (x=0, y=0)
\draw[waveBlue, very thick, ->, shorten >=2pt]
    (-1.0, \wfStartY)
    -- (\wfEndX, \wfEndY)
    node[above right, waveBlue] {Wavefront};

% ── Propagation direction arrow ──────────────────────────────────────────
\pgfmathsetmacro{\px}{sin(\theta)}
\pgfmathsetmacro{\py}{-cos(\theta)}
\draw[waveBlue, ->, thick]
    (3.0, 4.5)
    -- ({3.0 + 2*\px}, {4.5 + 2*\py})
    node[right, waveBlue, font=\sffamily\footnotesize]
    {Propagation direction};

% ── Bearing angle arc ────────────────────────────────────────────────────
% θ measured from broadside (vertical)
\draw[annoOrange, thin]
    (3.0, 4.5) -- (3.0, 6.5);   % broadside reference (vertical)
\draw[annoOrange, ->, thin]
    (3.0, 5.8) arc[start angle=90, end angle={90-\theta}, radius=1.3];
\node[annoOrange] at (3.6, 5.6) {$\theta$};

% ── Path-length difference construction ─────────────────────────────────
% For element 1 at x=d, y=0:
% The extra path length relative to element 0 is d*cos(theta).
% Draw right-triangle from element 0 to element 1 footprint on wavefront.

\pgfmathsetmacro{\footx}{\d * sin(\theta) * sin(\theta)}
\pgfmathsetmacro{\footy}{\d * sin(\theta) * (-cos(\theta))}
% Foot of perpendicular from element 1 onto the wavefront line through elem 0
% Wavefront unit normal: (sin(theta), -cos(theta))
% Projection of element-1 position (\d, 0) onto wavefront normal:
% proj = d * sin(theta)
% Foot = elem1 - proj * normal = (\d - d*sin^2, d*sin*cos)

\pgfmathsetmacro{\ex}{\d - \d*sin(\theta)*sin(\theta)}
\pgfmathsetmacro{\ey}{\d * sin(\theta) * cos(\theta)}

% Draw dashed path-length line from elem 1 down to wavefront
\draw[pathRed, dashed, thick] (\d, 0) -- (\ex, \ey);

% Path-length label
\pgfmathsetmacro{\midpx}{(\d + \ex)/2 + 0.3}
\pgfmathsetmacro{\midpy}{(\ey)/2}
\node[pathRed, right, font=\sffamily\footnotesize] at (\midpx, \midpy)
    {$d\cos\theta$};

% Right-angle mark at foot
\pgfmathsetmacro{\marksize}{0.12}
\draw[pathRed, thin]
    (\ex, \ey)
    -- ({\ex + \marksize*cos(\theta)}, {\ey + \marksize*sin(\theta)})
    -- ({\ex + \marksize*cos(\theta) - \marksize*sin(\theta)},
        {\ey + \marksize*sin(\theta) + \marksize*cos(\theta)})
    -- ({\ex - \marksize*sin(\theta)}, {\ey + \marksize*cos(\theta)});

% ── Arrival time timeline (below each element) ──────────────────────────
\def\timelineY{-1.4}
\def\timelineLen{0.9}

% Timeline axis
\draw[black!40, thin, ->]
    (-0.3, \timelineY) -- ({7*\d + 1.0}, \timelineY)
    node[right, font=\sffamily\footnotesize] {$t$};
\node[left, font=\sffamily\footnotesize] at (-0.3, \timelineY) {$\tau$};

% Dots for each sensor arrival time
% tau_n proportional to n; max tau at n=7 mapped to 0.7 units offset
\pgfmathsetmacro{\tscale}{0.5}   % scale: 1 unit of n maps to this many x-units
\foreach \i in {0,1,...,7}{
    \pgfmathsetmacro{\arrivalx}{\i * \d + \i * \tscale * 0.12}
    % vertical dashed line from element to timeline
    \draw[black!25, dashed, thin]
        ({\i*\d}, -0.3) -- ({\i*\d}, {\timelineY + 0.1});
    % dot
    \filldraw[sensorBlue]
        ({\i*\d + \i*0.08}, \timelineY) circle (0.07);
}
\node[font=\sffamily\tiny, text=black!60] at (3.5*\d, \timelineY - 0.4)
    {Progressive arrival: $\tau_n = n \cdot \dfrac{d\cos\theta}{c}$};

% ── Legend box ──────────────────────────────────────────────────────────
\draw[black!40, rounded corners=4pt, thin]
    (10.5, 1.5) rectangle (15.5, 5.5);
\node[above, font=\sffamily\footnotesize\bfseries] at (13.0, 5.2)
    {Legend};
\node[anchor=west, font=\sffamily\footnotesize] at (10.8, 4.6)
    {$\theta$ — bearing from broadside};
\node[anchor=west, font=\sffamily\footnotesize] at (10.8, 4.0)
    {$d = 0.1875\,\text{m}$ — element spacing};
\node[anchor=west, font=\sffamily\footnotesize] at (10.8, 3.4)
    {$c = 1500\,\text{m/s}$ — sound speed};
\node[anchor=west, font=\sffamily\footnotesize] at (10.8, 2.8)
    {$\tau_n = n \cdot \dfrac{d\cos\theta}{c}$ — delay at sensor $n$};
\node[anchor=west, font=\sffamily\footnotesize] at (10.8, 2.0)
    {$\mathbf{Y} = [\mathbf{y}_0\;\mathbf{y}_1\;\cdots\;\mathbf{y}_{N-1}]
     \in\mathbb{R}^{T\times N}$};

\end{tikzpicture}
\end{document}
